% Motor: XeLaTeX, fijado en .latexmkrc y en la receta del editor.
\documentclass[aspectratio=169,10pt,c]{beamer}
% Sin atajos de babel: así las comillas " del código se ven tal cual.
\PassOptionsToPackage{es-noshorthands}{babel}

% !TEX root = slides.tex
% Preambulo comun: estilo UTP, colores, fuentes, plantillas de beamer y
% comandos de apoyo didactico. Se carga desde slides.tex tras \documentclass.

% La opción "c" de beamer reparte el espacio libre 1:1.5 (arriba:abajo), por lo
% que el contenido queda por encima del centro. Igualamos ambos rellenos.
\makeatletter
\define@key{beamerframe}{c}[true]{%
  \beamer@frametopskip=0pt plus 1fill\relax%
  \beamer@framebottomskip=0pt plus 1fill\relax%
  \beamer@frametopskipautobreak=0pt plus .5\paperheight\relax%
  \beamer@framebottomskipautobreak=0pt plus .5\paperheight\relax%
}
\makeatother

\usepackage{graphicx}
\usepackage{tikz}
\usepackage{xcolor}
\usepackage{iftex}

\ifPDFTeX
  \usepackage[utf8]{inputenc}
  \usepackage[T1]{fontenc}
  \usepackage{lmodern}          % tipografia escalable para el respaldo con pdflatex
  \renewcommand{\familydefault}{\sfdefault}
\else
  \usepackage{fontspec}
  \IfFontExistsTF{Arial}
    {\setsansfont{Arial}[BoldFont={Arial Bold}]}
    {\setsansfont{Helvetica}[BoldFont={Helvetica Bold}]}
\fi

\usepackage[spanish,es-nodecimaldot]{babel}
\usepackage{microtype}
\usepackage{booktabs}
\usepackage{amssymb}
\usepackage{array}
\usetikzlibrary{arrows.meta,positioning,fit,calc,shapes.geometric}

\graphicspath{{assets/}}

% ---------- Identidad visual UTP ----------
\definecolor{utpRojo}{HTML}{C42B35}
\definecolor{utpNegro}{HTML}{111111}
\definecolor{utpGris}{HTML}{5F5F5F}
\definecolor{utpGrisMedio}{HTML}{9A9A9A}
\definecolor{utpGrisClaro}{HTML}{F1F1F1}
\definecolor{utpNaranja}{HTML}{C1670B}

\usetheme{default}
\setbeamercolor{normal text}{fg=utpNegro,bg=white}
\setbeamercolor{structure}{fg=utpRojo}
\setbeamercolor{frametitle}{fg=utpNegro,bg=white}
\setbeamercolor{title}{fg=utpNegro}
\setbeamercolor{subtitle}{fg=utpNegro}
\setbeamercolor{author}{fg=utpNegro}
\setbeamercolor{date}{fg=utpNegro}
\setbeamercolor{block title}{fg=white,bg=utpNegro}
\setbeamercolor{block body}{fg=utpNegro,bg=utpGrisClaro}
\setbeamercolor{alertblock title}{fg=white,bg=utpRojo}
\setbeamercolor{alertblock body}{fg=utpNegro,bg=utpRojo!7}
\setbeamercolor{exampleblock title}{fg=white,bg=utpGris}
\setbeamercolor{exampleblock body}{fg=utpNegro,bg=utpGris!8}

\setbeamerfont{title}{series=\bfseries,size=\fontsize{19.5}{22.5}\selectfont}
\setbeamerfont{subtitle}{size=\fontsize{11.5}{13.5}\selectfont}
\setbeamerfont{date}{series=\bfseries,size=\fontsize{11.5}{13.5}\selectfont}
\setbeamerfont{author}{size=\fontsize{10.5}{12.5}\selectfont}
\setbeamerfont{frametitle}{series=\bfseries,size=\large}
\setbeamerfont{block title}{series=\bfseries,size=\small}
\setbeamerfont{block body}{size=\normalsize}

\setbeamertemplate{navigation symbols}{}
\setbeamertemplate{headline}{}
\setbeamertemplate{blocks}[default]
\setbeamertemplate{itemize item}{\color{utpRojo}\small\raise.15ex\hbox{$\blacktriangleright$}}
\setbeamertemplate{itemize subitem}{\color{utpGris}\scriptsize$\bullet$}
\setbeamertemplate{enumerate item}{\color{utpRojo}\insertenumlabel.}
\setbeamertemplate{background canvas}{\color{white}\rule{\paperwidth}{\paperheight}}

% Título de diapositiva: negrita negra a la izquierda, logo UTP a la derecha
\setbeamertemplate{frametitle}{%
  \begin{tikzpicture}[remember picture,overlay]
    \node[anchor=north east,inner sep=0pt]
      at ([xshift=-7mm,yshift=-5.5mm]current page.north east)
      {\includegraphics[width=17mm]{utp-logo.png}};
  \end{tikzpicture}%
  \vspace{1.2mm}%
  \begin{minipage}[t]{.80\textwidth}\raggedright\insertframetitle\end{minipage}\par
  \vspace{1.1mm}{\color{utpRojo}\rule{\textwidth}{1.1pt}}
}

% Solo numeración discreta; la plantilla UTP no lleva pie de página
\setbeamertemplate{footline}{%
  \hfill{\color{utpGrisMedio}\scriptsize\insertframenumber}\hspace{6mm}\vspace{2.5mm}%
}

% Portada UTP
\setbeamertemplate{title page}{%
  \begin{tikzpicture}[remember picture,overlay]
    \node[anchor=center,inner sep=0pt] at (current page.center) {%
      \includegraphics[width=\paperwidth,height=\paperheight]{utp-cover-bg.jpeg}%
    };
    \node[anchor=north,align=center,text=utpNegro,inner sep=0pt,
      text width=145mm,font=\usebeamerfont{title}]
      at ([yshift=-25.2mm]current page.north) {\inserttitle};
    \node[anchor=north,align=center,text=utpNegro,inner sep=0pt,
      text width=142mm,font=\usebeamerfont{subtitle}]
      at ([yshift=-38.2mm]current page.north) {\insertsubtitle};
    \node[anchor=north,align=center,text=utpNegro,inner sep=0pt,
      font=\usebeamerfont{date}]
      at ([yshift=-50.4mm]current page.north) {\insertdate};
    \node[anchor=north,align=center,text=utpNegro,inner sep=0pt,
      font=\usebeamerfont{author}]
      at ([yshift=-58.4mm]current page.north) {\insertauthor};
  \end{tikzpicture}%
}

% ---------- Comandos de apoyo didáctico ----------
\newcommand{\concepto}[2]{%
  \begin{block}{#1}\large #2\end{block}}
\newcommand{\pregunta}[1]{%
  \begin{alertblock}{Decide y explica}#1\end{alertblock}}
\newcommand{\analogia}[1]{%
  \begin{exampleblock}{En palabras simples}\small #1\end{exampleblock}}
\newcommand{\enlarealidad}[1]{%
  \begin{block}{Cómo se ve en la realidad}\small #1\end{block}}
\newcommand{\ojo}[1]{%
  \begin{alertblock}{Ojo con esto}\small #1\end{alertblock}}
\newcommand{\apoyo}{{\normalfont\footnotesize\color{utpGris}~· material de apoyo}}
\newcommand{\marca}[1]{\tikz[baseline=(m.base)]{%
  \node[rounded corners=1.5pt,fill=utpGrisClaro,draw=utpGrisMedio,inner xsep=3pt,inner ysep=1.5pt,
        font=\footnotesize] (m) {#1};}}

\tikzset{
  caja/.style={draw=utpNegro, very thick, rounded corners=1.5pt, fill=white,
    align=center, minimum height=0.72cm, inner sep=5pt, font=\small},
  cajaazul/.style={caja, draw=utpGris, fill=white},
  cajaverde/.style={caja, draw=utpNegro, fill=utpGrisClaro},
  cajanaranja/.style={caja, draw=utpRojo!75, fill=utpRojo!7},
  flecha/.style={-{Stealth[length=2.3mm]}, very thick, draw=utpNegro},
  flechavuelta/.style={{Stealth[length=2.3mm]}-{Stealth[length=2.3mm]}, very thick, draw=utpNegro}
}

% ---------- Notación pata de gallo (crow's foot) ----------
% La rotación indica hacia dónde está la entidad: 0 derecha, 180 izquierda,
% 90 arriba, -90 abajo. La pata se abre hacia la entidad del lado "muchos".
\tikzset{
  relacion/.style={thick, draw=utpNegro},
  pics/uno/.style={code={%
    \draw[thick,utpNegro,line cap=round] (0,-1.5mm) -- (0,1.5mm);}},
  pics/muchos/.style={code={%
    \draw[thick,utpNegro,line cap=round,line join=round]
      (2.6mm,1.5mm) -- (0,0) -- (2.6mm,-1.5mm);
    \draw[thick,utpNegro,line cap=round] (0,0) -- (2.6mm,0);}},
}
\newcommand{\simbuno}{\tikz[baseline=-0.6ex]{\draw[relacion](0,0)--(7mm,0);\pic at (4.4mm,0){uno};}}
\newcommand{\simbmuchos}{\tikz[baseline=-0.6ex]{\draw[relacion](0,0)--(7mm,0);\pic at (4.4mm,0){muchos};}}

\usepackage{colortbl}
\usetikzlibrary{decorations.pathreplacing}

\title{Funciones y procedimientos almacenados}
\subtitle{Variables, \texttt{IF} y las piezas de un procedimiento en MySQL y SQL Server}
\date{Semana 14 - Sesión 28}
\author{Rel Guzman}

% Nombre de tabla, siempre en naranja
\newcommand{\tbl}[1]{\textcolor{utpNaranja}{\texttt{#1}}}
% Palabra reservada de SQL dentro de un bloque de código
\newcommand{\kw}[1]{\textcolor{utpRojo}{\textbf{#1}}}
% Comentario dentro de un bloque de código
\newcommand{\com}[1]{\textcolor{utpGrisMedio}{#1}}
% Texto de apoyo en gris
\newcommand{\gris}[1]{\textcolor{utpGris}{#1}}
% Insignia de práctica
\newcommand{\ins}[1]{\tikz[baseline=-0.5ex]{\node[rounded corners=2pt,fill=utpRojo,text=white,
  font=\bfseries\scriptsize,inner xsep=5pt,inner ysep=2.5pt] {#1};}}
% Etiqueta oscura: nombre de un operador o de un gestor
\newcommand{\etiqueta}[1]{\tikz[baseline=-0.5ex]{\node[rounded corners=2pt,fill=utpNegro,text=white,
  font=\bfseries\scriptsize,inner xsep=6pt,inner ysep=3pt] {#1};}}
% Número de paso
\newcommand{\npieza}[1]{\tikz[baseline=-0.6ex]{\node[circle,fill=utpRojo,text=white,
  font=\bfseries\scriptsize,inner sep=1.6pt,minimum size=4.2mm] {#1};}}
% Celda o texto resaltado en rojo claro
\newcommand{\mal}[1]{{\setlength{\fboxsep}{1.2pt}\colorbox{utpRojo!18}{#1}}}
% Lo importante, resaltado en amarillo dentro de un bloque de código
\newcommand{\dentro}[1]{{\setlength{\fboxsep}{1.5pt}\colorbox{yellow!45}{#1}}}
% Fila de una tabla resaltada
\newcommand{\fila}{\rowcolor{yellow!45}}
\newcommand{\filagris}{\rowcolor{utpGrisClaro}}
% Marcas de sí y no
\newcommand{\si}{$\checkmark$}
\newcommand{\no}{\textcolor{utpRojo}{$\times$}}
% Bloques de código sobre fondo gris: normal, chico y mínimo
\newsavebox{\sqlbuf}
\newenvironment{sqlcode}[1][.86\textwidth]
  {\begin{center}\begin{lrbox}{\sqlbuf}\begin{minipage}{#1}\ttfamily\footnotesize\vspace{6pt}\raggedright}
  {\vspace{6pt}\end{minipage}\end{lrbox}\colorbox{utpGrisClaro}{\usebox{\sqlbuf}}\end{center}}
\newenvironment{sqlchico}[1][.86\textwidth]
  {\begin{center}\begin{lrbox}{\sqlbuf}\begin{minipage}{#1}\ttfamily\scriptsize\vspace{5pt}\raggedright}
  {\vspace{5pt}\end{minipage}\end{lrbox}\colorbox{utpGrisClaro}{\usebox{\sqlbuf}}\end{center}}
\newenvironment{sqlmini}[1][.86\textwidth]
  {\begin{center}\begin{lrbox}{\sqlbuf}\begin{minipage}{#1}\ttfamily\fontsize{7}{8.3}\selectfont\vspace{4pt}\raggedright}
  {\vspace{4pt}\end{minipage}\end{lrbox}\colorbox{utpGrisClaro}{\usebox{\sqlbuf}}\end{center}}
\newenvironment{sqlmicro}[1][.86\textwidth]
  {\begin{center}\begin{lrbox}{\sqlbuf}\begin{minipage}{#1}\ttfamily\fontsize{6.4}{7.5}\selectfont\vspace{4pt}\raggedright}
  {\vspace{4pt}\end{minipage}\end{lrbox}\colorbox{utpGrisClaro}{\usebox{\sqlbuf}}\end{center}}
% Ancho de una columna menos el borde del bloque de código
\newcommand{\col}{\dimexpr\textwidth-2\fboxsep\relax}
% Sangría dentro de un bloque de código
\newcommand{\ind}{\hspace*{1.2em}}
% Mensaje de error dentro de un bloque de código
\newcommand{\err}[1]{\textcolor{utpRojo}{#1}}
% Anatomía de una instrucción: cada pieza del código con su pregunta al lado
\tikzset{
  pieza/.style={fill=utpGrisClaro,font=\ttfamily\footnotesize,anchor=west,inner sep=5pt,
    minimum height=0.62cm},
  para/.style={font=\small,anchor=west,text=utpNegro}
}
\newcommand{\piezax}{7.6}
\newcommand{\pieza}[3]{%
  \node[pieza] at (0,#1) {#2};
  \node[para] at (\piezax,#1) {\textcolor{utpRojo}{$\blacktriangleleft$}\ #3};}
% Diapositiva de sección: una palabra bien grande
\newcommand{\seccion}[2][1]{%
\begin{frame}[plain]
  \begin{tikzpicture}[remember picture,overlay]
    \node[anchor=north east,inner sep=0pt]
      at ([xshift=-10mm,yshift=-8mm]current page.north east)
      {\includegraphics[width=43mm]{utp-logo.png}};
    \node[anchor=center,inner sep=0pt] at (current page.center)
      {\resizebox{100mm}{!}{\bfseries\color{utpNegro}#2}};
    \ifnum#1=1
    \node[anchor=north,inner sep=0pt] at ([yshift=-11mm]current page.center)
      {\color{utpRojo}\rule{100mm}{1.6pt}};
    \fi
  \end{tikzpicture}
\end{frame}}
% Cabecera de columna dentro de una tabla en monoespaciado
\newcommand{\cab}[1]{\normalfont\bfseries #1}
% Resultado de una consulta: \res{columnas}{cabecera}{filas}
\newcommand{\res}[3]{%
  \centering{\scriptsize\color{utpGris}resultado}\\[3pt]
  \footnotesize\ttfamily
  \begin{tabular}{#1}
  \toprule
  #2 \\
  \midrule
  #3
  \bottomrule
  \end{tabular}}
% Lo mismo, cuando debajo sigue más contenido del frame
\newcommand{\resg}[3]{{\res{#1}{#2}{#3}\par}}
% Tabla pequeña con su nombre arriba: \mini{nombre}{columnas}{cabecera}{filas}
\newcommand{\mini}[4]{%
  \begin{tabular}{@{}c@{}}
  \tbl{#1}\\[3pt]
  {\scriptsize\ttfamily
  \begin{tabular}{#2}
  \toprule
  #3 \\
  \midrule
  #4
  \bottomrule
  \end{tabular}}
  \end{tabular}}
% Una tabla con título gris, alineada arriba: \tabla{título}{columnas}{cabecera}{filas}
\newcommand{\tabla}[4]{%
  \begin{tabular}[t]{@{}c@{}}
  {\scriptsize\color{utpGris}#1}\\[3pt]
  {\footnotesize\ttfamily
  \begin{tabular}[t]{#2}
  \toprule
  #3 \\
  \midrule
  #4
  \bottomrule
  \end{tabular}}
  \end{tabular}}
% La pregunta del frame, en una caja naranja centrada
\newcommand{\preg}[1]{{\centering\tikz\node[caja,cajanaranja,font=\normalsize,inner sep=6pt]{#1};\par}\vspace{4pt}}
% Forma de lo que devuelve una consulta: celdas con borde
\newcommand{\forma}[2]{{\ttfamily\footnotesize\renewcommand{\arraystretch}{1.25}%
  \begin{tabular}{#1}\hline #2\end{tabular}}}
% Estilos de los diagramas de esta sesión
\tikzset{
  codigo/.style={font=\ttfamily\footnotesize,inner sep=3pt},
  sub/.style={codigo,draw=utpRojo,very thick,fill=yellow!45,rounded corners=2pt},
  decision/.style={diamond,aspect=2.6,draw=utpRojo,very thick,fill=utpRojo!7,
    font=\scriptsize,align=center,inner sep=1.5pt},
  paso/.style={caja,font=\scriptsize,minimum height=.6cm,inner sep=4pt},
  % Una instrucción en un diagrama de pasos y lo que queda debajo
  inst/.style={draw=utpNegro,thick,rounded corners=1.5pt,fill=utpGrisClaro,font=\ttfamily\scriptsize,
    minimum height=.8cm,minimum width=2.2cm,align=center,inner sep=3pt},
  valor/.style={font=\footnotesize,align=center,text width=2.5cm},
  carril/.style={font=\small\bfseries,anchor=east}
}
% ---------------------------------------------------------------------
% Ventana de MySQL Workbench dibujada (sección de procedimientos)
\definecolor{wbBarra}{HTML}{E6E9ED}
% Íconos de Workbench
\newcommand{\rayo}{\tikz[baseline=-0.35ex]\fill[yellow!85!orange,draw=orange!70!black,line width=.3pt]
  (0.07,0.17)--(-0.06,-0.01)--(0.01,-0.01)--(-0.04,-0.17)--(0.09,0.04)--(0.02,0.04)--cycle;}
\newcommand{\rayocursor}{\rayo\kern-0.5pt{\fontsize{6}{6}\selectfont\bfseries I}}
\newcommand{\okwb}{\tikz[baseline={([yshift=-0.6ex]current bounding box.center)}]{\node[circle,fill=green!55!black,text=white,inner sep=0.4pt,
  font=\fontsize{5}{5}\selectfont]{$\checkmark$};}}
\newcommand{\errwb}{\tikz[baseline={([yshift=-0.6ex]current bounding box.center)}]{\node[circle,fill=utpRojo,text=white,inner sep=0.6pt,
  font=\fontsize{5}{5}\selectfont\bfseries]{$\times$};}}
\newcommand{\refrescar}{\tikz[baseline=-0.6ex]{\draw[-{Stealth[length=2pt]},thick,utpGris] (0.09,0) arc (0:290:0.09);}}
\newcommand{\carpeta}{\tikz[baseline=-0.4ex]{\fill[yellow!70!orange] (0,0) rectangle (0.24,0.15);
  \fill[yellow!70!orange] (0,0.15) rectangle (0.1,0.2);}}
\newcommand{\disco}{\tikz[baseline=-0.4ex]{\fill[blue!45!black] (0,0) rectangle (0.2,0.2);
  \fill[white] (0.05,0.12) rectangle (0.15,0.2);}}
\newcommand{\sqlmas}{\tikz[baseline=-0.6ex]{\node[draw=utpGris,fill=white,rounded corners=1pt,inner sep=1.5pt,
  font=\fontsize{5.5}{6}\selectfont\bfseries]{SQL\textcolor{green!50!black}{+}};}}
\newcommand{\abierto}{\textcolor{utpGris}{$\blacktriangledown$}\,}
\newcommand{\cerrado}{\textcolor{utpGris}{$\blacktriangleright$}\,}
\newcommand{\iconobase}{\tikz[baseline=-0.1ex]{\fill[cyan!45!black] (0,0) rectangle (0.15,0.15);}\,}
% Código del editor con número de línea: \lineas{1 & código \\ 2 & código \\}
\newcommand{\lineas}[1]{\begin{tabular}{@{}>{\color{utpGrisMedio}}r@{\hspace{8pt}}l@{}}#1\end{tabular}}
% Panel Output: \salida{icono & # & acción & mensaje \\}
\newcommand{\salida}[1]{{\fontsize{6.5}{8}\selectfont\begin{tabular}{@{}c@{\hspace{5pt}}c@{\hspace{8pt}}l@{\hspace{14pt}}l@{}}
  & \textcolor{utpGris}{\#} & \textcolor{utpGris}{Action} & \textcolor{utpGris}{Message} \\ #1\end{tabular}}}
% Result Grid: \grilla{columnas}{cabecera}{filas}
\newcommand{\grilla}[3]{{\ttfamily\scriptsize\setlength\tabcolsep{5pt}\renewcommand{\arraystretch}{1.15}%
  \begin{tabular}{#1}\hline\rowcolor{wbBarra}#2 \\\hline #3\hline\end{tabular}}}
% La ventana completa: \wbventana{árbol del Navigator}{código}{título del panel de abajo}{panel de abajo}
% Nombra: sqlmas, refrescar, arbol, pestana, rayo, rayoc, editor, paneltit, panel.
\newcommand{\wbventana}[4]{%
  \draw[utpGrisMedio,fill=white] (0,0) rectangle (12.4,-5.1);
  \fill[wbBarra] (0,0) rectangle (12.4,-0.45); \draw[utpGrisMedio] (0,-0.45)--(12.4,-0.45);
  \node (sqlmas) at (0.4,-0.225) {\sqlmas};
  \foreach \x in {1.0,1.45,1.9,2.35} \fill[utpGrisMedio!50] (\x-0.12,-0.11) rectangle (\x+0.12,-0.34);
  \fill[wbBarra] (0,-0.45) rectangle (3.2,-0.8);
  \node[anchor=west,font=\scriptsize\bfseries] at (0.08,-0.625) {Navigator};
  \node[anchor=west,font=\fontsize{5.5}{6}\selectfont\bfseries,text=utpGris] at (0.08,-1.0) {SCHEMAS};
  \node[inner sep=1pt] (refrescar) at (2.95,-1.0) {\refrescar};
  \node[anchor=north west,font=\scriptsize,align=left,inner sep=2pt] (arbol) at (0.08,-1.2) {#1};
  \draw[utpGrisMedio] (3.2,-0.45)--(3.2,-5.1);
  \node[anchor=south west,draw=utpGrisMedio,fill=white,font=\scriptsize,inner sep=2pt] (pestana) at (3.3,-0.8) {Query 1};
  \fill[wbBarra] (3.2,-0.8) rectangle (12.4,-1.22); \draw[utpGrisMedio] (3.2,-0.8)--(12.4,-0.8) (3.2,-1.22)--(12.4,-1.22);
  \node (abrir) at (3.5,-1.01) {\carpeta};
  \node (guardar) at (3.85,-1.01) {\disco};
  \node[inner sep=1.5pt] (rayo) at (4.3,-1.01) {\rayo};
  \node[inner sep=1.5pt] (rayoc) at (4.72,-1.01) {\rayocursor};
  \foreach \x in {5.2,5.55,5.9} \fill[utpGrisMedio!50] (\x-0.11,-0.9) rectangle (\x+0.11,-1.12);
  \node[anchor=north west,font=\ttfamily\scriptsize,align=left,inner sep=3pt] (editor) at (3.25,-1.27) {#2};
  \fill[wbBarra] (3.2,-3.35) rectangle (12.4,-3.7); \draw[utpGrisMedio] (3.2,-3.35)--(12.4,-3.35) (3.2,-3.7)--(12.4,-3.7);
  \node[anchor=west,font=\scriptsize\bfseries] (paneltit) at (3.3,-3.525) {#3};
  \node[anchor=north west,font=\scriptsize,inner sep=3pt] (panel) at (3.25,-3.75) {#4};
}
% Marca roja con número alrededor de un nodo de la ventana: \senal{nodo}{número}
\newcommand{\senal}[2]{%
  \draw[utpRojo,very thick,rounded corners=2pt] ([shift={(-2pt,-2pt)}]#1.south west) rectangle ([shift={(2pt,2pt)}]#1.north east);
  \node[circle,fill=utpRojo,text=white,font=\bfseries\scriptsize,inner sep=1.2pt,minimum size=4.2mm]
    at ([shift={(3pt,3pt)}]#1.north east) {#2};}
% Igual, con el número abajo a la derecha
\newcommand{\senalabajo}[2]{%
  \draw[utpRojo,very thick,rounded corners=2pt] ([shift={(-2pt,-2pt)}]#1.south west) rectangle ([shift={(2pt,2pt)}]#1.north east);
  \node[circle,fill=utpRojo,text=white,font=\bfseries\scriptsize,inner sep=1.2pt,minimum size=4.2mm]
    at ([shift={(3pt,-3pt)}]#1.south east) {#2};}
% Árbol del Navigator con la base cupones abierta
\newcommand{\arbolcupones}[1]{%
  \abierto\iconobase\textbf{cupones}\\
  \hspace*{8pt}\cerrado Tables\\
  \hspace*{8pt}\cerrado Views\\
  #1}
\begin{document}
{\setbeamertemplate{footline}{}\maketitle}
\begin{frame}[c]
\begin{block}{Logro de la sesión}
Al finalizar, podrás escribir \textbf{funciones} con \textbf{variables} e \textbf{IF} en MySQL y SQL Server, y reconocer las piezas de un \textbf{procedimiento almacenado}.
\end{block}
\end{frame}
% =======================================================================
% Revisión de la sesión anterior
% ---------------------------------------------------------------------
\begin{frame}{Revisión: ¿por qué la webapp usa SQL?}
\small \texttt{Sesion27/webapp} muestra la tabla \tbl{contacto}. Esto pasa cada vez que alguien entra a la página:
\vspace{0.35cm}
\begin{center}
\begin{tikzpicture}
  \tikzset{bd/.style={cylinder,shape border rotate=90,aspect=0.25,draw=utpRojo!75,very thick,fill=utpRojo!7,
                      minimum width=2cm,minimum height=1.5cm,align=center,font=\ttfamily\scriptsize},
           et/.style={font=\scriptsize,align=center}}
  % Navegador
  \node[draw=utpNegro,very thick,fill=white,minimum width=3.9cm,minimum height=2.3cm] (w) at (0,0) {};
  \node[anchor=north,fill=utpGrisClaro,draw=utpNegro,very thick,minimum width=3.9cm,minimum height=.42cm,font=\ttfamily\tiny] (b) at (w.north) {localhost:3000};
  \node[anchor=north west,font=\scriptsize\bfseries] at ([xshift=2mm,yshift=-.5mm]b.south west) {Agenda};
  \node[anchor=north,font=\tiny] at ([yshift=-5mm]b.south) {%
    \renewcommand{\arraystretch}{0.95}\setlength\tabcolsep{3pt}%
    \begin{tabular}{|l|l|l|}\hline
    \bfseries id & \bfseries nombre & \bfseries telefono \\\hline
    1 & Ana Torres & 987654321 \\\hline
    2 & Luis Ramos & 912345678 \\\hline
    3 & María Quispe & 956789123 \\\hline
    \end{tabular}};
  \node[font=\footnotesize\bfseries] at (0,1.5) {navegador};
  % app.js
  \node[caja,cajanaranja,minimum width=2.4cm,minimum height=1.3cm,font=\small,align=center] (app) at (4.9,0) {\texttt{app.js}\\[1pt]\scriptsize la página};
  % Base
  \node[bd] (db) at (10.6,0) {agenda.db};
  \node[font=\footnotesize\bfseries] at (10.6,1.5) {base de datos};
  % Flechas
  \draw[flecha] ([yshift=4mm]w.east) -- node[et,above]{\npieza{1}\\entras} ([yshift=4mm]app.west);
  \draw[flecha] ([yshift=-4mm]app.west) -- node[et,below]{tabla HTML\\\npieza{4}} ([yshift=-4mm]w.east);
  \draw[flecha] ([yshift=4mm]app.east) -- node[et,above]{\npieza{2}\\\texttt{\kw{SELECT} * \kw{FROM} contacto}} ([yshift=4mm]db.west);
  \draw[flecha] ([yshift=-4mm]db.west) -- node[et,below]{3 filas\\\npieza{3}} ([yshift=-4mm]app.east);
\end{tikzpicture}
\end{center}
\vspace{0.3cm}
\concepto{¿Por qué?}{\small Los contactos \textbf{no están escritos en la página}: están en la base. Si alguien agrega uno con \texttt{INSERT}, la página lo muestra al recargar, \textbf{sin cambiar} \texttt{app.js}.}
\end{frame}
% ---------------------------------------------------------------------
\begin{frame}{Revisión: cambia el SQL, cambia la página}
\small En \texttt{app.js} solo cambiamos el texto del SQL, guardamos y recargamos el navegador:
\vspace{0.2cm}
\begin{columns}[T,onlytextwidth]
\column{.47\textwidth}
\centering\etiqueta{Antes}
\begin{sqlmini}[\col]
.prepare("\kw{SELECT} * \kw{FROM} contacto")\\
\strut
\end{sqlmini}
\vspace{-4pt}
\begin{tikzpicture}
  \node[draw=utpNegro,very thick,fill=white,minimum width=4.6cm,minimum height=2.2cm] (w) at (0,0) {};
  \node[anchor=north,fill=utpGrisClaro,draw=utpNegro,very thick,minimum width=4.6cm,minimum height=.42cm,font=\ttfamily\tiny] (b) at (w.north) {localhost:3000};
  \node[anchor=north west,font=\scriptsize\bfseries] at ([xshift=2mm,yshift=-.5mm]b.south west) {Agenda};
  \node[anchor=north,font=\scriptsize] at ([yshift=-5mm]b.south) {%
    \renewcommand{\arraystretch}{0.95}\setlength\tabcolsep{3pt}%
    \begin{tabular}{|l|l|l|}\hline
    \bfseries id & \bfseries nombre & \bfseries telefono \\\hline
    1 & Ana Torres & 987654321 \\\hline
    2 & Luis Ramos & 912345678 \\\hline
    3 & María Quispe & 956789123 \\\hline
    \end{tabular}};
\end{tikzpicture}
\column{.47\textwidth}
\centering\etiqueta{Después}
\begin{sqlmini}[\col]
.prepare("\kw{SELECT} \dentro{nombre, telefono} \kw{FROM} contacto\\
\ind\ind\ind\dentro{\kw{ORDER BY} nombre \kw{DESC}}")
\end{sqlmini}
\vspace{-4pt}
\begin{tikzpicture}
  \node[draw=utpNegro,very thick,fill=white,minimum width=4.6cm,minimum height=2.2cm] (w) at (0,0) {};
  \node[anchor=north,fill=utpGrisClaro,draw=utpNegro,very thick,minimum width=4.6cm,minimum height=.42cm,font=\ttfamily\tiny] (b) at (w.north) {localhost:3000};
  \node[anchor=north west,font=\scriptsize\bfseries] at ([xshift=2mm,yshift=-.5mm]b.south west) {Agenda};
  \node[anchor=north,font=\scriptsize] at ([yshift=-5mm]b.south) {%
    \renewcommand{\arraystretch}{0.95}\setlength\tabcolsep{3pt}%
    \begin{tabular}{|l|l|}\hline
    \bfseries nombre & \bfseries telefono \\\hline
    María Quispe & 956789123 \\\hline
    Luis Ramos & 912345678 \\\hline
    Ana Torres & 987654321 \\\hline
    \end{tabular}};
\end{tikzpicture}
\end{columns}
\vspace{0.3cm}
\concepto{¿Por qué?}{\small La página muestra \textbf{exactamente} lo que devuelve la consulta: sus columnas, sus filas y su orden. Quien escribe el SQL decide qué ve el usuario.}
\end{frame}
% ---------------------------------------------------------------------
\begin{frame}{Revisión: el sistema de la feria por dentro}
\small \texttt{Sesion28/sistema\_implementado} (se abre con \texttt{npm start}). Cada acción en la pantalla termina en SQL:
\vspace{0.2cm}
\begin{center}
\begin{tikzpicture}
  \tikzset{pant/.style={caja,font=\scriptsize,minimum width=2.6cm,minimum height=.62cm},
           sql/.style={fill=utpGrisClaro,font=\ttfamily\scriptsize,anchor=west,inner sep=3pt,minimum height=.62cm},
           bd/.style={cylinder,shape border rotate=90,aspect=0.25,draw=utpRojo!75,very thick,fill=utpRojo!7,
                      minimum width=1.7cm,minimum height=1.4cm,align=center,font=\ttfamily\scriptsize}}
  \node[font=\footnotesize\bfseries] at (0,0.75) {En la pantalla};
  \node[font=\footnotesize\bfseries] at (6.1,0.75) {SQL en \texttt{src/db.js}};
  \foreach \y/\p/\q in {
      0/{Ver la entrada}/{\kw{SELECT} precio, disponibles\ldots{} \kw{FROM} entrada},
      -0.85/{Guardar precio}/{\kw{UPDATE} entrada \kw{SET} precio = ?, disponibles = ?\ldots},
      -1.7/{Ver las ventas}/{\kw{SELECT} \ldots{} \kw{FROM} venta \kw{ORDER BY} fecha\_venta \kw{DESC}}}{
    \node[pant] (p) at (0,\y) {\p};
    \node[sql] (q) at (2.1,\y) {\q};
    \draw[flecha] (p.east) -- (q.west);
  }
  \node[pant,cajanaranja,font=\scriptsize\bfseries] (pp) at (0,-2.55) {Pagar};
  \node[sql,fill=yellow!45] (qq) at (2.1,-2.55) {\kw{BEGIN} \ldots{} \kw{UPDATE} entrada + \kw{INSERT INTO} venta \ldots{} \kw{COMMIT}};
  \draw[flecha] (pp.east) -- (qq.west);
  \node[bd] (db) at (11.6,-1.27) {feria.db};
  \draw[utpGrisMedio,thick,decorate,decoration={brace,amplitude=5pt}] (10.2,0.3) -- (10.2,-2.85);
  \draw[flecha] (10.35,-1.27) -- (db.west);
\end{tikzpicture}
\end{center}
\vspace{0.2cm}
\concepto{¿Por qué?}{\small Los datos de la feria viven en la base. La aplicación solo \textbf{traduce cada clic a SQL}, y todo ese SQL está en un solo archivo: \texttt{src/db.js}.}
\end{frame}
% ---------------------------------------------------------------------
\begin{frame}{Revisión: ¿por qué la feria usa una transacción?}
\small En \texttt{Sesion28/sistema\_implementado}, vender una entrada son \textbf{dos cambios}. Quedan 499 entradas y hay 1 venta (500 en total). ¿Y si el \texttt{INSERT} falla?
\vspace{0.15cm}
\begin{center}
\begin{tikzpicture}
  \tikzset{inst/.append style={minimum width=2cm},
           fin/.style={caja,font=\scriptsize,text width=2.6cm,align=center,minimum height=1cm}}
  % Sin transacción
  \node[anchor=west] at (-1,0.8) {\etiqueta{Sin transacción}};
  \node[inst] (a1) at (0,0) {\kw{UPDATE} entrada\\disponibles - 1};
  \node[inst] (a2) at (2.5,0) {\kw{INSERT INTO}\\venta \ldots};
  \node[fin,draw=utpRojo,fill=utpRojo!7] (a3) at (10.3,0) {498 + 1 venta = \textbf{499}\\\textbf{se perdió una entrada}};
  \draw[flecha] (a1)--(a2); \draw[flecha] (a2)--(a3);
  \node[valor] at (0,-.72) {\gris{498, ya guardado}};
  \node[valor] at (2.5,-.72) {\err{\textbf{falla}}};
  % Con transacción
  \node[anchor=west] at (-1,-1.55) {\etiqueta{Con transacción: lo que hace src/db.js}};
  \node[inst,fill=utpNegro,text=white] (b0) at (0,-2.35) {\textbf{BEGIN}\\\textbf{TRANSACTION}};
  \node[inst] (b1) at (2.5,-2.35) {\kw{UPDATE} entrada\\disponibles - 1};
  \node[inst] (b2) at (5,-2.35) {\kw{INSERT INTO}\\venta \ldots};
  \node[inst,fill=utpRojo!7,draw=utpRojo!75] (b3) at (7.5,-2.35) {\kw{ROLLBACK};};
  \node[fin,fill=yellow!45] (b4) at (10.3,-2.35) {499 + 1 venta = \textbf{500}\\\textbf{no se pierde nada}};
  \draw[flecha] (b0)--(b1); \draw[flecha] (b1)--(b2); \draw[flecha] (b2)--(b3); \draw[flecha] (b3)--(b4);
  \node[valor] at (2.5,-3.07) {\gris{498, pendiente}};
  \node[valor] at (5,-3.07) {\err{\textbf{falla}}};
  \node[valor] at (7.5,-3.07) {\gris{vuelve a 499}};
\end{tikzpicture}
\end{center}
\vspace{0.2cm}
\concepto{¿Por qué?}{\small Cada entrada descontada debe tener \textbf{su venta}. Por eso \texttt{src/db.js} hace los dos cambios en una transacción: si todo sale bien, \texttt{COMMIT}; si algo falla, \texttt{ROLLBACK} deshace los dos.}
\end{frame}
% ---------------------------------------------------------------------
\begin{frame}{Revisión: la compra paso a paso}
\small Al pagar, \texttt{registrarCompra} revisa todo \textbf{dentro} de la transacción. Si algo no se cumple, \texttt{ROLLBACK}:
\vspace{0.3cm}
\begin{center}
\begin{tikzpicture}
  \tikzset{chk/.style={caja,font=\scriptsize,text width=1.45cm,align=center,minimum height=1.05cm,inner sep=3pt},
           no/.style={caja,draw=utpRojo,fill=utpRojo!7,font=\tiny,text width=1.85cm,align=center,minimum height=.9cm,inner sep=2pt}}
  \node[inst,fill=utpNegro,text=white,minimum width=1.6cm] (b) at (0,0) {\textbf{BEGIN}\\\textbf{TRANSACTION}};
  \node[chk] (c1) at (2.3,0)  {¿venta abierta?};
  \node[chk] (c2) at (4.6,0)  {¿DNI nuevo?};
  \node[chk] (c3) at (6.9,0)  {¿pago nuevo?};
  \node[chk] (c4) at (9.2,0)  {\texttt{\kw{UPDATE}} -1\\¿quedaba alguna?};
  \node[chk,fill=yellow!45] (ok) at (11.5,0) {\texttt{\kw{INSERT}}\\venta\\\texttt{\textbf{COMMIT}}};
  \draw[flecha] (b)--(c1);
  \foreach \a/\z in {c1/c2,c2/c3,c3/c4,c4/ok} \draw[flecha] (\a) -- node[above,font=\tiny\bfseries]{sí} (\z);
  \node[no] (n1) at (2.3,-1.75) {\texttt{\textbf{ROLLBACK}}\\«La venta de entradas está cerrada.»};
  \node[no,very thick] (n2) at (4.6,-1.75) {\texttt{\textbf{ROLLBACK}}\\«Este DNI ya tiene una entrada comprada\ldots»};
  \node[no] (n3) at (6.9,-1.75) {\texttt{\textbf{ROLLBACK}}\\«Este pago ya fue registrado.»};
  \node[no] (n4) at (9.2,-1.75) {\texttt{\textbf{ROLLBACK}}\\«Ya no quedan entradas disponibles.»};
  \foreach \i in {1,2,3,4} \draw[flecha] (c\i) -- node[right,font=\tiny\bfseries]{no} (n\i);
  \node[font=\scriptsize\bfseries,text=utpRojo] at (4.6,-2.6) {Luis otra vez};
\end{tikzpicture}
\end{center}
\vspace{0.2cm}
\concepto{¿Por qué?}{\small Las preguntas y los cambios van en la \textbf{misma} transacción. Así, una compra rechazada \textbf{no deja nada a medias}: ni una entrada descontada sin venta, ni una venta sin entrada.}
\end{frame}
% =======================================================================
\seccion{Programación en SQL}
\begin{frame}{Caso de estudio: sistema de cupones de descuento}
\begin{columns}[T,onlytextwidth]
\column{.47\textwidth}
\small Una tienda en línea lanza \textbf{códigos de descuento}.
\vspace{0.2cm}
\begin{itemize}\setlength\itemsep{2pt}
  \item El \textbf{responsable de promociones} registra los cupones y los activa o desactiva.
  \item El \textbf{cliente} escribe el código al pagar.
  \item Un cupón descuenta un \textbf{porcentaje} o un \textbf{monto fijo} en soles.
  \item Tiene fecha de inicio y, si se quiere, fecha de fin.
\end{itemize}
\column{.5\textwidth}
\centering
\begin{tikzpicture}
  \tikzset{panel/.style={draw=utpNegro,very thick,fill=white,minimum width=6.4cm,minimum height=3.3cm,anchor=north}}
  \node[panel] (p) at (0,0) {};
  \node[anchor=north,fill=utpNegro,text=white,minimum width=6.4cm,minimum height=.55cm,font=\footnotesize\bfseries]
    (b) at (p.north) {Tienda: pagar mi compra};
  \node[anchor=north west,font=\small] at ([xshift=4mm,yshift=-3mm]b.south west) {Total de la compra};
  \node[anchor=north east,font=\small\bfseries] at ([xshift=-4mm,yshift=-3mm]b.south east) {S/ 120.00};
  \node[anchor=north west,font=\small] at ([xshift=4mm,yshift=-10mm]b.south west) {Cupón};
  \node[anchor=north east,draw=utpGrisMedio,font=\ttfamily\small,minimum width=2.8cm,inner sep=3pt]
    at ([xshift=-4mm,yshift=-9.4mm]b.south east) {BIENVENIDA10};
  \draw[utpGrisMedio] ([xshift=4mm,yshift=-17mm]b.south west) -- ([xshift=-4mm,yshift=-17mm]b.south east);
  \node[anchor=north west,font=\small\bfseries] at ([xshift=4mm,yshift=-19mm]b.south west) {A pagar};
  \node[anchor=north east,font=\small\bfseries,text=utpRojo] at ([xshift=-4mm,yshift=-19mm]b.south east) {S/ ?};
\end{tikzpicture}
\end{columns}
\vspace{0cm}
\begin{center}
\mini{cupon}{rllrlll}{\cab{id} & \cab{codigo} & \cab{tipo\_descuento} & \cab{valor\_descuento} & \cab{estado} & \cab{vigente\_desde} & \cab{vigente\_hasta}}{%
1 & BIENVENIDA10 & porcentaje & 10.00 & activo & 2026-01-01 & NULL \\
2 & NOVIEMBRE20 & monto\_fijo & 20.00 & activo & 2026-11-01 & 2026-12-01 \\
3 & VERANO15 & porcentaje & 15.00 & activo & 2026-01-01 & 2026-04-01 \\
4 & NAVIDAD25 & porcentaje & 25.00 & activo & 2026-12-15 & 2026-12-26 \\
5 & PRUEBA50 & monto\_fijo & 50.00 & inactivo & 2026-10-01 & NULL \\}
\end{center}
\end{frame}
% ---------------------------------------------------------------------
\begin{frame}{¿Cuánto paga el cliente? La fórmula va en la base}
\small Una compra de \textbf{S/ 120}. El cálculo depende del tipo de cupón:
\vspace{0.15cm}
\begin{center}
\small
\renewcommand{\arraystretch}{1.3}
\begin{tabular}{@{}llll@{}}
\toprule
\bfseries Cupón & \bfseries Tipo & \bfseries Cálculo & \bfseries Paga \\
\midrule
\texttt{BIENVENIDA10} & porcentaje: 10 & $120 - 120 \times 10 / 100$ & \textbf{S/ 108.00} \\
\texttt{NOVIEMBRE20} & monto fijo: 20 & $120 - 20$ & \textbf{S/ 100.00} \\
\bottomrule
\end{tabular}
\end{center}
\vspace{0.15cm}
\begin{columns}[c,onlytextwidth]
\column{.5\textwidth}
\centering
\begin{tikzpicture}
  \foreach \t/\y/\i in {Panel de cupones/0/1, Tienda en línea/-0.8/2}
    \node[paso,minimum width=2.6cm] (a\i) at (0,\y) {\t};
  \node[caja,cajanaranja,font=\scriptsize,align=center] (f) at (3.6,-0.4) {base de datos\\[2pt]\texttt{precio\_final()}};
  \foreach \i in {1,2} \draw[flecha] (a\i.east) -- (f.west);
\end{tikzpicture}
\column{.46\textwidth}
\raggedright\small Los \textbf{dos sistemas} necesitan el mismo cálculo. Si lo escribe cada uno, el panel puede mostrar un precio y la tienda cobrar otro.\\[6pt]
Se escribe \textbf{una vez}, en la base, como \textbf{función}, y todos la usan.
\end{columns}
\vspace{0.2cm}
\concepto{Función definida por el usuario}{\small Una función que \textbf{tú escribes} con \texttt{CREATE FUNCTION}. Queda \textbf{guardada en la base} y se usa en un \texttt{SELECT}, igual que \texttt{ROUND} o \texttt{UPPER}.}
\end{frame}
% ---------------------------------------------------------------------
\begin{frame}{Una función en MySQL y en SQL Server}
\small El caso más simple: un cupón de \textbf{porcentaje}.
\begin{columns}[T,onlytextwidth]
\column{.47\textwidth}
\centering\etiqueta{MySQL}
\begin{sqlmini}[\col]
\kw{CREATE FUNCTION} aplicar\_porcentaje(\\
\ind\mal{p\_precio} \kw{DECIMAL}(10,2),\\
\ind\mal{p\_porcentaje} \kw{DECIMAL}(5,2))\\
\kw{RETURNS DECIMAL}(10,2)\\
\kw{DETERMINISTIC}\\
\kw{RETURN} p\_precio - p\_precio * p\_porcentaje / 100;\\[6pt]
\kw{SELECT} aplicar\_porcentaje(120, 10);\ \com{-- 108.00}
\end{sqlmini}
\column{.5\textwidth}
\centering\etiqueta{SQL Server}
\begin{sqlmini}[\col]
\kw{CREATE} \mal{\kw{OR ALTER}} \kw{FUNCTION} \mal{dbo.}aplicar\_porcentaje(\\
\ind\mal{@precio} \kw{DECIMAL}(10,2),\\
\ind\mal{@porcentaje} \kw{DECIMAL}(5,2))\\
\kw{RETURNS DECIMAL}(10,2)\\
\mal{\kw{AS}}\\
\mal{\kw{BEGIN}}\\
\ind\kw{RETURN} @precio - @precio * @porcentaje / 100;\\
\mal{\kw{END}};\\
\mal{\kw{GO}}\\[4pt]
\kw{SELECT} \mal{dbo.}aplicar\_porcentaje(120, 10);\ \com{-- 108.00}
\end{sqlmini}
\end{columns}
\vspace{0.1cm}
\begin{center}
\footnotesize
\renewcommand{\arraystretch}{1.1}
\begin{tabular}{@{}>{\bfseries}lll@{}}
\toprule
 & \bfseries MySQL & \bfseries SQL Server \\
\midrule
Parámetros & \texttt{p\_precio} (\texttt{p\_}: costumbre) & \texttt{@precio}: la \texttt{@} es obligatoria \\
Cuerpo & \texttt{RETURN \ldots;} & entre \texttt{AS BEGIN} y \texttt{END} \\
Para usarla & \texttt{aplicar\_porcentaje(\ldots)} & \texttt{dbo.aplicar\_porcentaje(\ldots)} \\
\bottomrule
\end{tabular}
\end{center}
\end{frame}
% ---------------------------------------------------------------------
\begin{frame}{Usarla sobre las filas de una tabla}
\small En el panel: con una compra de S/ 120, ¿cuánto pagaría el cliente con cada cupón de porcentaje?
\begin{columns}[c,onlytextwidth]
\column{.55\textwidth}
\begin{sqlcode}[\col]
\kw{SELECT} codigo, valor\_descuento,\\
\ind\dentro{aplicar\_porcentaje(120, valor\_descuento)}\\
\ind\ind\kw{AS} a\_pagar\\
\kw{FROM} cupon\\
\kw{WHERE} tipo\_descuento = 'porcentaje'\\
\kw{ORDER BY} id;
\end{sqlcode}
\column{.42\textwidth}
\res{lrr}{\cab{codigo} & \cab{valor\_descuento} & \cab{a\_pagar}}{BIENVENIDA10 & 10.00 & 108.00 \\ VERANO15 & 15.00 & 102.00 \\ NAVIDAD25 & 25.00 & 90.00 \\}
\end{columns}
\vspace{0.3cm}
\concepto{Idea clave}{\small La función se ejecuta \textbf{una vez por cada fila}: cada vez recibe el \texttt{valor\_descuento} de esa fila. En SQL Server se escribe \texttt{dbo.aplicar\_porcentaje}.}
\end{frame}
% ---------------------------------------------------------------------
\begin{frame}{Un cupón de monto fijo necesita variables e \texttt{IF}}
\small \texttt{NOVIEMBRE20} descuenta \textbf{S/ 20}, no el 20 \%. Con la función de antes, el cálculo sale mal:
\begin{columns}[c,onlytextwidth]
\column{.55\textwidth}
\begin{sqlcode}[\col]
\kw{SELECT} aplicar\_porcentaje(120, 20);
\end{sqlcode}
\column{.4\textwidth}
\centering
{\scriptsize\color{utpGris}resultado}\\[2pt]
\mal{\ttfamily 96.00}\quad\small debía ser \textbf{100.00}
\end{columns}
\vspace{0.3cm}
\small La función \texttt{precio\_final(precio, cupon)} tiene que hacer \textbf{tres pasos}:
\vspace{0.3cm}
\begin{center}
\begin{tikzpicture}
  \node[paso,text width=3.6cm] (a) at (0,0) {\npieza{1} Leer el cupón y \textbf{guardar} su tipo y su valor};
  \node[paso,text width=3.6cm] (b) at (4.6,0) {\npieza{2} \textbf{Según el tipo}, calcular el descuento};
  \node[paso,text width=3.6cm] (c) at (9.2,0) {\npieza{3} Devolver el precio menos el descuento};
  \draw[flecha] (a)--(b); \draw[flecha] (b)--(c);
  \node[font=\small\bfseries,text=utpRojo] at (0,-1) {variables};
  \node[font=\small\bfseries,text=utpRojo] at (4.6,-1) {\texttt{IF}};
  \node[font=\small\bfseries,text=utpRojo] at (9.2,-1) {\texttt{RETURN}};
\end{tikzpicture}
\end{center}
\end{frame}
% ---------------------------------------------------------------------
\begin{frame}{Variables: guardar un valor para usarlo después}
\concepto{Variable}{\small Un \textbf{nombre} que guarda \textbf{un valor} mientras se ejecuta tu código. Se le da un valor con \texttt{SET}.}
\vspace{0.1cm}
\begin{columns}[T,onlytextwidth]
\column{.48\textwidth}
\centering\etiqueta{MySQL}
\begin{sqlchico}[\col]
\kw{SET} @precio = 120;\\
\kw{SET} @descuento = 20;\\
\kw{SELECT} @precio - @descuento \kw{AS} a\_pagar;\\
\com{-- 100}
\end{sqlchico}
\footnotesize Guardar lo que devuelve un \texttt{SELECT}:
\begin{sqlchico}[\col]
\kw{SELECT} valor\_descuento \dentro{\kw{INTO} @descuento}\\
\kw{FROM} cupon \kw{WHERE} codigo = 'NOVIEMBRE20';
\end{sqlchico}
\column{.48\textwidth}
\centering\etiqueta{SQL Server}
\begin{sqlchico}[\col]
\kw{DECLARE} @precio \kw{DECIMAL}(10,2) = 120;\\
\kw{DECLARE} @descuento \kw{DECIMAL}(10,2) = 20;\\
\kw{SELECT} @precio - @descuento \kw{AS} a\_pagar;\\
\com{-- 100.00}
\end{sqlchico}
\footnotesize Guardar lo que devuelve un \texttt{SELECT}:
\begin{sqlchico}[\col]
\kw{SELECT} \dentro{@descuento = valor\_descuento}\\
\kw{FROM} cupon \kw{WHERE} codigo = 'NOVIEMBRE20';
\end{sqlchico}
\end{columns}
\ojo{Dentro de una función de MySQL, las variables se crean con \texttt{DECLARE} y van \textbf{sin} \texttt{@}: \texttt{DECLARE v\_tipo VARCHAR(10);}}
\end{frame}
% ---------------------------------------------------------------------
\begin{frame}{\texttt{IF}: elegir un camino}
\begin{columns}[c,onlytextwidth]
\column{.42\textwidth}
\centering
\begin{tikzpicture}
  \node[decision] (d) at (0,0) {¿el tipo es\\'porcentaje'?};
  \node[paso,text width=2.4cm] (si) at (-1.65,-1.9) {descuento =\\precio $\times$ valor / 100};
  \node[paso,text width=2.1cm] (no) at (1.65,-1.9) {descuento =\\valor};
  \draw[flecha] (d.west) -| node[above,pos=.25,font=\scriptsize\bfseries]{sí} (si.north);
  \draw[flecha] (d.east) -| node[above,pos=.25,font=\scriptsize\bfseries]{no} (no.north);
\end{tikzpicture}\\[8pt]
\small Se ejecuta \textbf{solo uno} de los dos caminos.
\column{.55\textwidth}
\centering\etiqueta{MySQL}
\begin{sqlcode}[\col]
\kw{IF} v\_tipo = 'porcentaje' \mal{\kw{THEN}}\\
\ind\kw{SET} v\_descuento = p\_precio * v\_valor / 100;\\
\kw{ELSE}\\
\ind\kw{SET} v\_descuento = v\_valor;\\
\mal{\kw{END IF}};
\end{sqlcode}
\centering\etiqueta{SQL Server}
\begin{sqlcode}[\col]
\kw{IF} @tipo = 'porcentaje'\\
\ind\kw{SET} @descuento = @precio * @valor / 100;\\
\kw{ELSE}\\
\ind\kw{SET} @descuento = @valor;
\end{sqlcode}
\end{columns}
\analogia{En SQL Server no hay \texttt{THEN} ni \texttt{END IF}. Si un camino tiene \textbf{varias} instrucciones, van entre \texttt{BEGIN} y \texttt{END}.}
\end{frame}
% ---------------------------------------------------------------------
\begin{frame}{Todo junto: \texttt{precio\_final} en MySQL}
\begin{columns}[c,onlytextwidth]
\column{.58\textwidth}
\begin{sqlmicro}[\col]
\dentro{\kw{DELIMITER} \$\$}\\
\kw{CREATE FUNCTION} precio\_final(p\_precio \kw{DECIMAL}(10,2),\\
\ind\ind\ind\ind\ind\ind\ind\ind\ind\ind p\_cupon\_id \kw{INT})\\
\kw{RETURNS DECIMAL}(10,2)\\
\kw{READS SQL DATA}\\
\kw{BEGIN}\\
\ind\kw{DECLARE} v\_tipo \kw{VARCHAR}(10);\\
\ind\kw{DECLARE} v\_valor \kw{DECIMAL}(10,2);\\
\ind\kw{DECLARE} v\_descuento \kw{DECIMAL}(10,2);\\[1.5pt]
\ind\kw{SELECT} tipo\_descuento, valor\_descuento \kw{INTO} v\_tipo, v\_valor\\
\ind\kw{FROM} cupon \kw{WHERE} id = p\_cupon\_id;\hfill\npieza{1}\\[1.5pt]
\ind\kw{IF} v\_tipo = 'porcentaje' \kw{THEN}\\
\ind\ind\kw{SET} v\_descuento = p\_precio * v\_valor / 100;\\
\ind\kw{ELSE}\\
\ind\ind\kw{SET} v\_descuento = v\_valor;\\
\ind\kw{END IF};\hfill\npieza{2}\\[1.5pt]
\ind\kw{IF} v\_descuento > p\_precio \kw{THEN}\\
\ind\ind\kw{SET} v\_descuento = p\_precio;\\
\ind\kw{END IF};\\[1.5pt]
\ind\kw{RETURN} p\_precio - v\_descuento;\hfill\npieza{3}\\
\kw{END}\dentro{\$\$}\\
\dentro{\kw{DELIMITER} ;}
\end{sqlmicro}
\column{.39\textwidth}
\raggedright\footnotesize
\textbf{\texttt{BEGIN} \ldots{} \texttt{END}}: agrupa varias instrucciones, cada una con su \texttt{;}.\\[5pt]
\textbf{\texttt{DELIMITER \$\$}}: le dice a MySQL que la función termina en \texttt{\$\$}, no en el primer \texttt{;}.\\[5pt]
\textbf{\texttt{READS SQL DATA}}: la función lee una tabla.\\[5pt]
El segundo \texttt{IF} evita que se pague menos de 0.
\vspace{0.2cm}
\begin{sqlmini}[\col]
\kw{SELECT} precio\_final(120, 1),\ \com{-- 108.00}\\
\ind\ind\ind precio\_final(120, 2),\ \com{-- 100.00}\\
\ind\ind\ind precio\_final(30, 5);\ \ \ \com{-- \ \ 0.00}
\end{sqlmini}
{\scriptsize\gris{Cupones 1, 2 y 5. Con \texttt{PRUEBA50} el descuento (50) es mayor que el precio (30): paga 0.}}
\end{columns}
\end{frame}
% ---------------------------------------------------------------------
\begin{frame}{\texttt{IF} con varios caminos: \texttt{ELSEIF}}
\renewcommand{\arraystretch}{0.85}
\small ¿En qué situación está cada cupón? Se revisa \textbf{en orden}: gana la primera condición que se cumple.
\begin{columns}[T,onlytextwidth]
\column{.5\textwidth}
\centering\etiqueta{MySQL}
\begin{sqlmini}[\col]
\kw{IF} v\_estado = 'inactivo' \kw{THEN}\\
\ind\kw{SET} v\_situacion = 'inactivo';\\
\dentro{\kw{ELSEIF}} p\_fecha < v\_desde \kw{THEN}\\
\ind\kw{SET} v\_situacion = 'por empezar';\\
\dentro{\kw{ELSEIF}} v\_hasta \kw{IS NOT NULL}\\
\ind\ind\kw{AND} p\_fecha >= v\_hasta \kw{THEN}\\
\ind\kw{SET} v\_situacion = 'vencido';\\
\kw{ELSE}\\
\ind\kw{SET} v\_situacion = 'vigente';\\
\kw{END IF};
\end{sqlmini}
\centering\etiqueta{SQL Server}
\begin{sqlmini}[\col]
\kw{IF} @estado = 'inactivo' \ldots\\
\dentro{\kw{ELSE IF}} @fecha < @desde \ldots\\
\dentro{\kw{ELSE IF}} @hasta \kw{IS NOT NULL AND} \ldots\\
\kw{ELSE} \ldots
\end{sqlmini}
\column{.47\textwidth}
\begin{sqlmini}[\col]
\kw{SELECT} codigo,\\
\ind situacion\_cupon(id, '2026-11-10') \kw{AS} situacion\\
\kw{FROM} cupon \kw{ORDER BY} id;
\end{sqlmini}
\res{ll}{\cab{codigo} & \cab{situacion}}{BIENVENIDA10 & vigente \\ NOVIEMBRE20 & vigente \\ VERANO15 & vencido \\ NAVIDAD25 & por empezar \\ PRUEBA50 & inactivo \\}
\par\vspace{4pt}
{\centering\normalfont\scriptsize\color{utpGris}La fecha es fija para que todos veamos lo mismo.\\En la aplicación se usa la fecha de hoy.\par}
\end{columns}
\end{frame}
% ---------------------------------------------------------------------
\begin{frame}{\ins{Práctica 1} La base de cupones y sus funciones}
\npieza{1}\ \small Crea la base \texttt{cupones}: dos scripts, en orden. Están en \texttt{ejemplos/practicas/}:
\vspace{0.1cm}
\begin{center}
\footnotesize
\renewcommand{\arraystretch}{1.3}
\setlength\tabcolsep{4pt}
\begin{tabular}{@{}lll@{}}
\toprule
\bfseries Gestor & \bfseries 1. Modelo físico (tablas) & \bfseries 2. Registros (datos) \\
\midrule
MySQL (Workbench) & \texttt{mysql/practica\_1\_modelo\_fisico.sql} & \texttt{mysql/practica\_1\_registros.sql} \\
SQL Server (SSMS) & \texttt{sqlserver/practica\_1\_modelo\_fisico.sql} & \texttt{sqlserver/practica\_1\_registros.sql} \\
\bottomrule
\end{tabular}
\end{center}
\vspace{0.1cm}
\npieza{2}\ \small Crea las funciones y pruébalas. Los scripts están en \texttt{ejemplos/cupones/}, en este orden:
\vspace{0.05cm}
\begin{center}
\footnotesize
\renewcommand{\arraystretch}{1.2}
\begin{tabular}{@{}lll@{}}
\toprule
\bfseries Gestor & \bfseries 1. Primera función & \bfseries 2. Variables e \texttt{IF} \\
\midrule
MySQL & \texttt{1\_funciones\_mysql.sql} & \texttt{2\_variables\_if\_mysql.sql} \\
SQL Server & \texttt{1\_funciones\_sqlserver.sql} & \texttt{2\_variables\_if\_sqlserver.sql} \\
\bottomrule
\end{tabular}
\end{center}
\vspace{0.05cm}
\npieza{3}\ \small Comprueba que al final del segundo script salgan solo los cupones vigentes: \textbf{BIENVENIDA10} paga \textbf{108.00} y \textbf{NOVIEMBRE20} paga \textbf{100.00}.
\end{frame}
% ---------------------------------------------------------------------
% =======================================================================
\seccion{Procedimientos almacenados}
\begin{frame}{Desactivar un cupón no es un cálculo}
\small El responsable de promociones quiere \textbf{desactivar} \texttt{VERANO15}. No hay nada que calcular: hay que \textbf{cambiar} una fila.
\vspace{0.3cm}
\begin{center}
\begin{tikzpicture}
  \tikzset{q/.style={caja,text width=5.3cm,minimum height=2.3cm,font=\small,align=left,anchor=north}}
  \node[q] (f) at (0,0) {\textbf{Función}\\[4pt]
    \footnotesize Recibe datos, \textbf{calcula} y \textbf{devuelve un valor}.\\[4pt]
    Se usa \textbf{dentro} de un \texttt{SELECT}:\\
    {\ttfamily\scriptsize \kw{SELECT} precio\_final(120, 1);}};
  \node[q,cajanaranja] (p) at (6.2,0) {\textbf{Procedimiento almacenado}\\[4pt]
    \footnotesize \textbf{Hace algo}: cambia datos (\texttt{UPDATE}, \texttt{INSERT}\ldots) o muestra un resultado.\\[4pt]
    Se ejecuta \textbf{solo}, con \texttt{CALL}:\\
    {\ttfamily\scriptsize \kw{CALL} desactivar\_cupon('VERANO15');}};
\end{tikzpicture}
\end{center}
\vspace{0.3cm}
\concepto{Procedimiento almacenado}{\small Un grupo de instrucciones con \textbf{nombre}, guardado \textbf{en la base}. Lo creas una vez y después lo usas cuantas veces quieras con \texttt{CALL}.}
\end{frame}
% ---------------------------------------------------------------------
\begin{frame}{Las piezas de un procedimiento}
\small Este procedimiento recibe un código y desactiva ese cupón.
\vspace{0.35cm}
\begin{center}
\renewcommand{\piezax}{6.1}
\begin{tikzpicture}
  \pieza{0}{\kw{CREATE PROCEDURE}}{crea algo con nombre y lo \textbf{guarda en la base}}
  \pieza{-0.85}{desactivar\_cupon}{su \textbf{nombre}: con él lo vas a usar}
  \pieza{-1.7}{(p\_codigo \kw{VARCHAR}(255))}{el \textbf{parámetro}: el dato que recibe}
  \pieza{-2.55}{\kw{UPDATE} cupon \kw{SET} estado = 'inactivo'}{lo que \textbf{hace} cada vez que lo usas}
  \pieza{-3.4}{\kw{WHERE} codigo = p\_codigo;}{usa el parámetro para elegir \textbf{la fila}}
\end{tikzpicture}
\end{center}
\vspace{0.3cm}
{\centering\footnotesize\gris{\texttt{p\_} es solo una costumbre: así no se confunde el parámetro \texttt{p\_codigo} con la columna \texttt{codigo}.}\par}
\end{frame}
% ---------------------------------------------------------------------
% ---------------------------------------------------------------------
\begin{frame}[plain]
  \begin{tikzpicture}[remember picture,overlay]
    \node[anchor=north west,font=\bfseries\fontsize{24}{28}\selectfont]
      at ([xshift=12mm,yshift=-9mm]current page.north west) {Preguntas};
    \node[anchor=north east,inner sep=0pt]
      at ([xshift=-10mm,yshift=-8mm]current page.north east)
      {\includegraphics[width=43mm]{utp-logo.png}};
    \node[anchor=center,inner sep=0pt]
      at ([yshift=-5mm]current page.center)
      {\includegraphics[width=78mm]{preguntas-ilustracion.png}};
  \end{tikzpicture}
\end{frame}
\begin{frame}{Cierre}
\begin{center}
\small
\renewcommand{\arraystretch}{1.12}
\begin{tabular}{>{\raggedright\bfseries}p{3.6cm}>{\raggedright\arraybackslash}p{9.4cm}}
\toprule
Función del usuario & Tu propio cálculo, guardado en la base con \texttt{CREATE FUNCTION}. Se usa en un \texttt{SELECT}. \\
Variable & Un nombre que guarda un valor: \texttt{DECLARE}, \texttt{SET}, \texttt{SELECT \ldots{} INTO}. \\
\texttt{IF} & Elige un camino. MySQL: \texttt{ELSEIF} y \texttt{END IF}; SQL Server: \texttt{ELSE IF}. \\
Procedimiento almacenado & Instrucciones con nombre, guardadas con \texttt{CREATE PROCEDURE}. No calcula: \textbf{hace} algo, como cambiar datos. Se usa con \texttt{CALL}. \\
\bottomrule
\end{tabular}
\end{center}
\end{frame}

\begin{frame}{Referencias}
\footnotesize
\begin{itemize}\setlength\itemsep{7pt}
\item Silberschatz, A., Korth, H. F., \& Sudarshan, S. (2014). \textit{Fundamentos de bases de datos} (6.ª ed.). McGraw-Hill.
\item Oracle. MySQL Reference Manual: \href{https://dev.mysql.com/doc/refman/8.0/en/create-procedure.html}{\textit{CREATE PROCEDURE and CREATE FUNCTION}}, \href{https://dev.mysql.com/doc/refman/8.0/en/call.html}{\textit{CALL Statement}} y \href{https://dev.mysql.com/doc/refman/8.0/en/if.html}{\textit{IF Statement}}.
\item Microsoft. Transact-SQL: \href{https://learn.microsoft.com/sql/t-sql/statements/create-function-transact-sql}{\textit{CREATE FUNCTION}}, \href{https://learn.microsoft.com/sql/t-sql/statements/create-procedure-transact-sql}{\textit{CREATE PROCEDURE}} y \href{https://learn.microsoft.com/sql/t-sql/language-elements/if-else-transact-sql}{\textit{IF...ELSE}}.
\end{itemize}
\end{frame}
% Cierre UTP
{
  \setbeamertemplate{background canvas}{%
    \includegraphics[width=\paperwidth,height=\paperheight]{gracias-final.pdf}%
  }
  \setbeamertemplate{footline}{}
  \begin{frame}[plain]
  \end{frame}
}

\end{document}
